\question[6]
Per ciascuno degli specificatori \lstinline|%x|, \lstinline|%p|, \lstinline|%s|, \lstinline|%n| indica cosa fa e la rispettiva nota di sicurezza. Spiega inoltre a cosa servono la \emph{larghezza} (\lstinline|%Nx|) e l'\emph{accesso diretto} (\lstinline|%n$|).
\begin{solution}
\begin{center}\small
\begin{tabular}{@{}lll@{}}
\toprule
Specificatore & Cosa fa & Nota di sicurezza\\
\midrule
\lstinline|%x| & stampa una parola in esadecimale & dump dello stack/registri\\
\lstinline|%p| & stampa un puntatore (indirizzo)   & rivela indirizzi $\to$ leak\\
\lstinline|%s| & stampa la stringa all'indirizzo fornito & dereferenzia $\to$ crash o leak\\
\lstinline|%n| & \emph{scrive} quanti caratteri stampati finora & primitiva di \textbf{scrittura}\\
\bottomrule
\end{tabular}
\end{center}
\textbf{Larghezza} \lstinline|%Nx|: forza l'emissione di almeno \lstinline|N| caratteri; il conteggio dei caratteri emessi è esattamente ciò che poi \lstinline|%n| scrive in memoria.\\
\textbf{Accesso diretto} \lstinline|%n$| (indicizzazione POSIX): \lstinline|%7$x| legge direttamente il 7\textsuperscript{o} argomento senza ripetere sette \lstinline|%x|. Si combina con la larghezza (es.\ \lstinline|%10$n|) per leggere/scrivere una posizione precisa.
\end{solution}
